%%%%%%%%%%%%%%%%%%%%%%%%%%%%%%%%%%%%%%%%%%%%%%%%%%%%%%%%%%%%%%%%%%%%%%%%%%%%%%
%
% 10_operads.tex
%
%%%%%%%%%%%%%%%%%%%%%%%%%%%%%%%%%%%%%%%%%%%%%%%%%%%%%%%%%%%%%%%%%%%%%%%%%%%%%%

\begin{mosbox}{10. Operadic Computation and Parallel Scheduling}

\BodyFont

Within MOS, a reasoning act is represented as a directed acyclic graph
(DAG) of operator nodes rather than a linear sequence of computations.

Each node corresponds to an operator, while directed edges encode
dependency relations. Independent computations therefore become explicit,
allowing reasoning to be executed in parallel whenever dependencies
permit.

\end{mosbox}

%%%%%%%%%%%%%%%%%%%%%%%%%%%%%%%%%%%%%%%%%%%%%%%%%%%%%%%%%%%%%%%%%%%%%%%%%%%%%%
%% HERO
%%%%%%%%%%%%%%%%%%%%%%%%%%%%%%%%%%%%%%%%%%%%%%%%%%%%%%%%%%%%%%%%%%%%%%%%%%%%%%

\begin{eqbox}

\[
\Huge
\boxed{
T=(N,\prec)
}
\]

\vspace{2mm}

\BodyFont

A reasoning program is a partially ordered set of operator nodes whose
execution order is determined by the dependency relation
$\prec$.

\end{eqbox}

%%%%%%%%%%%%%%%%%%%%%%%%%%%%%%%%%%%%%%%%%%%%%%%%%%%%%%%%%%%%%%%%%%%%%%%%%%%%%%
%% OPERAD TREE
%%%%%%%%%%%%%%%%%%%%%%%%%%%%%%%%%%%%%%%%%%%%%%%%%%%%%%%%%%%%%%%%%%%%%%%%%%%%%%

\begin{center}

\begin{tikzpicture}[
  scale=1.2,
  op/.style={circle, draw=MOSGold, line width=2pt, fill=MOSRed, text=white, font=\Large\bfseries, minimum size=10mm},
  leaf/.style={circle, draw=MOSRed, line width=1.5pt, text=white, font=\Large\bfseries, minimum size=8mm},
  edge from parent/.style={draw=MOSRed, line width=2pt},
  level 1/.style={sibling distance=2.5cm, level distance=2cm},
  level 2/.style={sibling distance=1.5cm, level distance=2cm},
  grow'=down
]

\node[op] (f) {$f$}
  child { node[op] (g) {$g$}
    child { node[leaf] (a) {$a$} }
    child { node[leaf] (b) {$b$} }
  }
  child { node[op] (h) {$h$} }
  child { node[op] (k) {$k$}
    child { node[leaf] (c) {$c$} }
  };

\end{tikzpicture}

\end{center}

%%%%%%%%%%%%%%%%%%%%%%%%%%%%%%%%%%%%%%%%%%%%%%%%%%%%%%%%%%%%%%%%%%%%%%%%%%%%%%
%% FOLIATION
%%%%%%%%%%%%%%%%%%%%%%%%%%%%%%%%%%%%%%%%%%%%%%%%%%%%%%%%%%%%%%%%%%%%%%%%%%%%%%

\begin{highlightbox}

\SectionTitle{Operadic Foliation}

\BodyFont

The scheduler partitions the reasoning DAG into successive
\emph{antichain layers}.

\[
L_0,\;
L_1,\;
L_2,\;
\ldots
\]

where

\[
L_0
=
\{\text{minimal nodes}\}
\]

and

\[
L_{\ell+1}
=
\{
n:
\text{all predecessors belong to }
L_{\le\ell}
\}.
\]

Each layer is executed in parallel because no two nodes within the same
antichain possess a dependency relation.

\end{highlightbox}

%%%%%%%%%%%%%%%%%%%%%%%%%%%%%%%%%%%%%%%%%%%%%%%%%%%%%%%%%%%%%%%%%%%%%%%%%%%%%%
%% SUPPORT
%%%%%%%%%%%%%%%%%%%%%%%%%%%%%%%%%%%%%%%%%%%%%%%%%%%%%%%%%%%%%%%%%%%%%%%%%%%%%%

\begin{mosbox}{Support Sets}

\BodyFont

Each operator records the subset of fine-complex vertices that it reads
and writes.

Operators with disjoint support sets may safely execute concurrently,
providing a mathematically justified scheduling rule rather than a
heuristic approximation.

\end{mosbox}

%%%%%%%%%%%%%%%%%%%%%%%%%%%%%%%%%%%%%%%%%%%%%%%%%%%%%%%%%%%%%%%%%%%%%%%%%%%%%%
%% TAKEAWAY
%%%%%%%%%%%%%%%%%%%%%%%%%%%%%%%%%%%%%%%%%%%%%%%%%%%%%%%%%%%%%%%%%%%%%%%%%%%%%%

\begin{remarkbox}

\BodyFont

\textbf{Interpretation}

Reasoning is viewed as composition rather than sequential execution.

The partial order determines correctness, while the foliation determines
parallel execution.

\end{remarkbox}